\ThesisFigureSlide{2.1}{The Waterfall model, used as a historical reference point for common software artifact roles}{28}{85 324 85 122}

\ThesisFigureSlide{3.1}{Table of contents from the expanded Smith et al. template}{51}{95 145 95 115}
\ThesisFigureSlide{3.2}{Table of units section from the GlassBR SRS}{52}{75 300 75 310}
\ThesisFigureSlide{3.3}{Table of symbols (truncated) section from the GlassBR SRS}{53}{75 430 75 95}
\ThesisFigureSlide{3.4}{Abbreviations and acronyms (truncated) section from the GlassBR SRS}{53}{75 115 75 400}
\ThesisFigureSlide{3.5}{Table of contents for the GlassBR Module Guide}{56}{80 150 80 160}
\ThesisFigureSlide{3.6}{Calc module from the GlassBR Module Guide}{57}{75 550 75 120}
\ThesisFigureSlide{3.7}{Python source code directory structure for GlassBR}{58}{100 330 100 110}
\ThesisFigureSlide{3.8}{Source code of the \texttt{Calc.py} module for GlassBR}{60}{70 90 70 100}
\ThesisFigureSlide{3.9}{Table of symbols (truncated) section from the GamePhysics SRS}{62}{80 430 80 110}
\ThesisFigureSlide{3.10}{Data definition for dimensionless load (\qhat{}) from the GlassBR SRS}{63}{80 455 80 105}
\ThesisFigureSlide{3.11}{Theoretical model of conservation of thermal energy found in both the SWHS and NoPCM SRSs}{65}{80 460 80 115}
\ThesisFigureSlide{3.12}{Instance model difference between SWHS and NoPCM}{66}{100 110 100 95}

\ThesisFigureSlide{4.1}{\texttt{NamedChunk} definition}{75}{80 526 80 117}
\ThesisFigureSlide{4.2}{Some example instances of \texttt{ConceptChunk} using the \texttt{dcc} smart constructor}{75}{80 371 80 302}
\ThesisFigureSlide{4.3}{Projecting knowledge into a chunk's definition}{77}{80 603 80 110}
\ThesisFigureSlide{4.4}{The force \texttt{ConceptChunk}}{77}{80 535 80 230}
\ThesisFigureSlide{4.5}{The force and acceleration \texttt{UnitalChunk}s}{77}{80 86 80 666}
\ThesisFigureSlide{4.6}{The fundamental SI units as captured in Drasil}{78}{80 562 80 122}
\ThesisFigureSlide{4.7}{Examples of chunks for units derived from other SI unit chunks}{78}{80 423 80 261}
\ThesisFigureSlide{4.8}{Defining real number multiplication in \texttt{Expr}}{80}{80 569 80 155}
\ThesisFigureSlide{4.9}{Encoding Newton's second law of motion}{80}{80 479 80 252}
\ThesisFigureSlide{4.10}{Newton's second law of motion as a \emph{chunk}}{81}{80 553 80 119}
\ThesisFigureSlide{4.11}{A representative subset of Drasil's chunk classification system showing some of the encoded property relationships}{83}{80 378 80 90}
\ThesisFigureSlide{4.12}{A representative subset of Drasil's chunks showing how chunks may belong to multiple classifications based on the underlying knowledge they project}{83}{80 140 80 509}
\ThesisFigureSlide{4.13}{Recipe language for an SRS}{85}{80 479 80 122}
\ThesisFigureSlide{4.14}{\texttt{ReqrmntSec} decomposition and definitions}{85}{80 306 80 349}
\ThesisFigureSlide{4.15}{SRS recipe for GlassBR}{86}{80 139 80 90}
\ThesisFigureSlide{4.16}{System information for GlassBR}{87}{80 479 80 117}
\ThesisFigureSlide{4.17}{The \texttt{iMods} definition}{89}{80 437 80 105}
\ThesisFigureSlide{4.18}{Recipe for GlassBR's executable code}{91}{80 468 80 112}
\ThesisFigureSlide{4.19}{Recipe language for executable code (\texttt{CodeSpec})}{92}{80 306 80 95}
\ThesisFigureSlide{4.20}{The \texttt{readTableMod} module definition}{93}{80 500 80 180}
\ThesisFigureSlide{4.21}{The \texttt{Choices} object definition}{93}{80 144 80 425}
\ThesisFigureSlide{4.22}{\texttt{genTeX} function for \LaTeX{} rendering}{97}{80 618 80 119}
\ThesisFigureSlide{4.23}{\texttt{generateCode} and \texttt{genPackage} for code rendering}{98}{80 488 80 117}
\ThesisFigureSlide{4.24}{\texttt{prntDoc'} function for document output}{98}{80 374 80 338}
\ThesisFigureSlide{4.25}{\texttt{prntMake} function for document output}{98}{80 99 80 621}
\ThesisFigureSlide{4.26}{\texttt{createCodeFiles} for code output}{99}{80 555 80 105}
\ThesisFigureSlide{4.27}{Definition of \qhat{} as a chunk in the GlassBR knowledge base}{100}{80 536 80 117}
\ThesisFigureSlide{4.28}{Raw generated \LaTeX{} for the \qhat{} data definition}{101}{80 148 80 95}
\ThesisFigureSlide{4.29}{\qhat{} as defined in generated Python code}{103}{80 580 80 117}
